\section{Cluster Operations}
\label{sec:cluster}

This section describes the \raft{}-facing control plane: the compact control
state, the commands that mutate it, the \graft{} state machine adapter, the
read queries, and the \texttt{vannak-node} binary that hosts them.

\subsection{ClusterControlState}

\texttt{ClusterControlState} is a dependency-free reducer for the compact state a
\raft{} state machine applies after log entries commit. Raw event traffic is
never routed through it. It holds:

\begin{description}
\item[\texttt{nodes}] the set of cluster member node ids.
\item[\texttt{placement\_maps}] epoch-keyed placement maps, bounded to five.
\item[\texttt{writer\_leases}] the current lease per \ipto{} target.
\item[\texttt{outbox\_checkpoints}] last acknowledged offsets per (shard, target).
\item[\texttt{sealed\_segments}] sealed segment manifests by id.
\item[\texttt{checkpoints}] recovery checkpoint manifests by id.
\end{description}

\subsection{ClusterControlCommand}

All mutations flow through a single command enum, so the set of replicated
state transitions is small and intentional.

\begin{lstlisting}
pub enum ClusterControlCommand {
    AddNode(NodeId),
    RemoveNode(NodeId),
    SetIptoPlacementMap(IptoPlacementMap),
    GrantWriterLease(WriterLease),
    RecordOutboxCheckpoint(MetadataOutboxCheckpoint),
    RecordSealedSegment(SegmentManifest),
    RecordCheckpoint(CheckpointManifest),
}
\end{lstlisting}

\subsection{State Machine Adapter}

The \texttt{raft} feature provides \texttt{ClusterStateMachine}, which wraps
\texttt{ClusterControlState} behind an \texttt{RwLock} so it can be shared
across \raft{} threads behind an \texttt{Arc}. It implements
\texttt{graft\_core::StateMachine} and \texttt{QueryableStateMachine}. Commands
and snapshots are encoded as JSON via dedicated serializable mirror types so the
core control types stay free of serialization concerns.

\subsection{Apply, Snapshot, and Restore Flow}

The \raft{} runtime drives the adapter through four operations:

\begin{description}
\item[\texttt{apply}] decodes the committed command and applies it under the
  write lock; a decode error is ignored so a single bad entry cannot crash the
  state machine.
\item[\texttt{apply\_with\_result}] like \texttt{apply} but returns an encoded
  error string when the command is rejected.
\item[\texttt{snapshot}] serializes the entire control state to JSON under the
  read lock.
\item[\texttt{restore}] replaces the state by replaying a decoded snapshot.
\end{description}

Restore is implemented by replaying each element of the snapshot through
\texttt{apply}, which means snapshot installation goes through the same
validation as live commands.

\subsection{Read Queries}

Read-only access uses \texttt{ClusterQuery}, served by the queryable state
machine without taking the write lock.

\begin{lstlisting}
pub enum ClusterQuery {
    ResolveTarget { shard_id: u64 },
    ResolveWithFallback { shard_id: u64 },
    ListIptoInstances,
    CurrentPlacementEpoch,
    GetWriterLease { target: String },
    GetOutboxCheckpoint { shard_id: u64, target: String },
    GetSealedSegment { segment_id: String },
}
\end{lstlisting}

Each query returns a JSON-encoded result, for example the resolved target name,
the ordered fallback candidate list, the current placement epoch, or a lease
holder.

\subsection{Node Lifecycle}

\texttt{AddNode} inserts a node into the membership set. \texttt{RemoveNode}
removes it and also drops every writer lease that node held, so authority does
not linger on a departed node. Membership is the basis for validating lease
grants.

\subsection{Placement Map Lifecycle}

\texttt{SetIptoPlacementMap} installs a new map. The proposed epoch must be
strictly greater than the current one or it is rejected with
\texttt{StalePlacementEpoch}. After insertion the history is pruned to the five
most recent epochs, supporting the query fallback described in
Section~\ref{sec:topology}.

\subsection{Writer Lease Lifecycle}

\texttt{GrantWriterLease} requires the holder to be a known node
(\texttt{UnknownLeaseHolder} otherwise) and a strictly greater
\texttt{LeaseEpoch} than any current lease for the target
(\texttt{StaleLeaseEpoch} otherwise). This gives single-writer authority per
\ipto{} target with monotonic fencing.

\subsection{Checkpoint Lifecycle}

\texttt{RecordOutboxCheckpoint} and \texttt{RecordCheckpoint} both enforce
monotonic epochs, rejecting stale records with \texttt{StaleCheckpointEpoch} and
\texttt{StaleCheckpointManifest} respectively. \texttt{RecordSealedSegment}
records a sealed \texttt{SegmentManifest} for the durable segment index.

\subsection{Error Variants}

\begin{table}[htbp]
\centering
\begin{tabular}{ll}
\hline
Variant & Condition \\
\hline
\texttt{NoPlacementSlots} & placement map built with no slots \\
\texttt{ZeroVnodes} & a slot has zero virtual nodes \\
\texttt{InvalidPlacementRange} & range end below start \\
\texttt{OverlappingPlacementRanges} & two overrides overlap \\
\texttt{StalePlacementEpoch} & placement epoch not advanced \\
\texttt{UnknownLeaseHolder} & lease holder is not a node \\
\texttt{StaleLeaseEpoch} & lease epoch not advanced \\
\texttt{StaleCheckpointEpoch} & outbox checkpoint epoch not advanced \\
\texttt{StaleCheckpointManifest} & checkpoint manifest epoch not advanced \\
\hline
\end{tabular}
\caption{\texttt{ClusterControlError} variants.}
\end{table}

\subsection{The vannak-node Binary}

The \texttt{node} feature builds \texttt{vannak-node}, which hosts a \raft{}
node backed by \texttt{ClusterStateMachine}. It uses \graft{} components: a
\texttt{FileLogStore}, a \texttt{FilePersistentStateStore}, a
\texttt{RaftRuntime}, and a \texttt{RaftHandler} dispatching over a TCP listener
with the \graft{} envelope codec. Peers are discovered from DNS SRV records or
an explicit peer list.

\begin{lstlisting}[numbers=none]
vannak-node serve <host> <port> <peer-id> <data-dir> <peers...>
vannak-node --srv _vannak._tcp.vannak.local <host> <port> <peer-id> <data-dir>
vannak-node probe <host> <port>
\end{lstlisting}

\begin{vnote}
On Linux the binary binds dual-stack (\texttt{[::]}) when given a wildcard or
loopback host, so a single listener accepts both IPv4 and IPv6 traffic inside
the container network.
\end{vnote}

\endinput
